\documentclass[11pt,a4paper]{article}
\usepackage[utf8]{inputenc}
\usepackage[T1]{fontenc}
\usepackage{graphicx}
\usepackage{amsmath}
\usepackage{amssymb}
\usepackage{booktabs}
\usepackage{hyperref}
\usepackage{float}
\usepackage{caption}
\usepackage{subcaption}
\usepackage[margin=2.5cm]{geometry}
\usepackage{natbib}
\usepackage{xcolor}

\title{\textbf{Impact of Pretraining Source and Transfer Learning Strategy on CNN Explainability:\\A Comparative Study Using Grad-CAM}}

\author{
    Fatemeh Elahe\\
    \texttt{[your-email@university.edu]}
}

\date{January 2026}

\begin{document}

\maketitle

% ============================================================================
% ABSTRACT
% ============================================================================
\begin{abstract}
Transfer learning has become a cornerstone of modern deep learning, yet the relationship between pretraining source, fine-tuning strategy, and model interpretability remains underexplored. This study investigates how different combinations of pretraining datasets (ImageNet vs. domain-specific Flowers-102) and transfer learning strategies (Linear Probing vs. Fine-Tuning) affect both classification performance and the explainability of Convolutional Neural Networks. We conduct a systematic 2×2 factorial experiment using ResNet-18 architecture on a 10-class ImageNet subset, evaluating four model configurations. Our results demonstrate that while Fine-Tuning consistently outperforms Linear Probing in accuracy, the pretraining source significantly influences where models focus their attention, as revealed through Grad-CAM visualizations. Models pretrained on ImageNet exhibit more semantically meaningful attention patterns compared to those pretrained on the domain-specific Flowers-102 dataset. These findings have important implications for practitioners choosing transfer learning strategies in applications where model interpretability is crucial.

\vspace{0.3cm}
\noindent\textbf{GitHub Repository:} \url{https://github.com/[username]/block-4-5-model-ab-fatemeh}
\end{abstract}

% ============================================================================
% 1. INTRODUCTION
% ============================================================================
\section{Introduction}

\subsection{Motivation and Background}

Deep learning has revolutionized computer vision, with Convolutional Neural Networks (CNNs) achieving remarkable performance on image classification tasks \citep{krizhevsky2012imagenet}. However, training deep networks from scratch requires vast amounts of labeled data and computational resources. Transfer learning addresses this challenge by leveraging knowledge from models pretrained on large datasets, enabling effective learning even with limited task-specific data \citep{yosinski2014transferable}.

The success of transfer learning raises fundamental questions about how pretrained representations influence model behavior. While performance metrics like accuracy are commonly reported, understanding \textit{why} a model makes certain predictions is equally important, particularly in high-stakes applications such as medical diagnosis, autonomous driving, and security systems \citep{arrieta2020explainable}.

Explainable AI (XAI) techniques, such as Gradient-weighted Class Activation Mapping (Grad-CAM) \citep{selvaraju2017grad}, provide visual explanations of CNN decisions by highlighting image regions that most influence predictions. However, the interplay between transfer learning choices and the quality of these explanations remains poorly understood.

\subsection{Related Work}

\textbf{Transfer Learning in Computer Vision.} \citet{yosinski2014transferable} demonstrated that features learned in early CNN layers are general and transferable, while later layers become increasingly task-specific. \citet{kornblith2019better} showed that models with higher ImageNet accuracy generally transfer better to other tasks, though this relationship is not always linear.

\textbf{Linear Probing vs. Fine-Tuning.} Two dominant strategies exist for transfer learning: (1) \textit{Linear Probing}, where only the final classification layer is trained while keeping the pretrained backbone frozen, and (2) \textit{Fine-Tuning}, where all layers are updated. \citet{kumar2022fine} found that linear probing can sometimes outperform fine-tuning when distribution shift exists between pretraining and target domains.

\textbf{Explainability in Transfer Learning.} While Grad-CAM \citep{selvaraju2017grad} and similar techniques have been widely adopted for explaining CNN predictions, few studies have systematically examined how transfer learning strategies affect explanation quality. \citet{adebayo2018sanity} highlighted that some saliency methods may not faithfully represent model reasoning, emphasizing the need for careful interpretation.

\subsection{Contribution}

This work makes the following contributions:

\begin{enumerate}
    \item We conduct a systematic 2×2 factorial experiment comparing pretraining sources (general-purpose ImageNet vs. domain-specific Flowers-102) and transfer strategies (Linear Probing vs. Fine-Tuning).

    \item We evaluate not only classification accuracy but also model explainability through comprehensive Grad-CAM analysis, examining how different configurations affect where models focus attention.

    \item We provide practical recommendations for practitioners selecting transfer learning approaches when model interpretability is a priority.
\end{enumerate}

% ============================================================================
% 2. METHODS
% ============================================================================
\section{Methods}

\subsection{Overview}

Our experimental framework evaluates four model configurations in a 2×2 factorial design, as illustrated in Table \ref{tab:experimental_design}. Each configuration uses ResNet-18 \citep{he2016deep} as the base architecture, differing only in pretraining source and transfer strategy.

\begin{table}[H]
\centering
\caption{Experimental Design: 2×2 Factorial Configuration}
\label{tab:experimental_design}
\begin{tabular}{lcc}
\toprule
& \textbf{Linear Probing} & \textbf{Fine-Tuning} \\
\midrule
\textbf{ImageNet Pretrained} & Model A & Model B \\
\textbf{Flowers-102 Pretrained} & Model C & Model D \\
\bottomrule
\end{tabular}
\end{table}

\subsection{Model Architecture}

We employ ResNet-18, a widely-used CNN architecture with approximately 11 million parameters. The architecture consists of:

\begin{itemize}
    \item \textbf{Backbone}: Four residual blocks (conv2\_x through conv5\_x) with skip connections that enable effective gradient flow during training.
    \item \textbf{Global Average Pooling}: Reduces spatial dimensions to a 512-dimensional feature vector.
    \item \textbf{Classification Head}: A fully-connected layer mapping the 512-dimensional features to the number of target classes.
\end{itemize}

\subsection{Pretraining Sources}

\textbf{ImageNet-1K} \citep{deng2009imagenet}: A large-scale dataset containing 1.2 million images across 1,000 diverse categories including animals, objects, and scenes. Models pretrained on ImageNet learn general-purpose visual features applicable to a wide range of downstream tasks.

\textbf{Flowers-102} \citep{nilsback2008automated}: A domain-specific dataset containing 8,189 images of 102 flower species. Models pretrained on this dataset learn features specialized for fine-grained botanical classification, potentially capturing texture and color patterns specific to flowers.

\subsection{Transfer Learning Strategies}

\textbf{Linear Probing (Models A, C):} The pretrained backbone is frozen (all parameters have \texttt{requires\_grad=False}), and only the final fully-connected layer is trained. This approach:
\begin{itemize}
    \item Preserves the original feature representations
    \item Requires fewer computational resources
    \item Uses a higher learning rate ($\eta = 10^{-3}$) since only the classifier is optimized
\end{itemize}

\textbf{Fine-Tuning (Models B, D):} All layers are trainable, allowing the model to adapt its feature representations to the target task. This approach:
\begin{itemize}
    \item Enables task-specific feature adaptation
    \item Requires more computational resources
    \item Uses a lower learning rate ($\eta = 10^{-4}$) to prevent catastrophic forgetting
\end{itemize}

\subsection{Explainability Method: Grad-CAM}

Gradient-weighted Class Activation Mapping (Grad-CAM) \citep{selvaraju2017grad} generates visual explanations by computing the gradient of the target class score with respect to the feature maps of the final convolutional layer. For a target class $c$, the importance weight $\alpha_k^c$ for feature map $k$ is:

\begin{equation}
\alpha_k^c = \frac{1}{Z} \sum_i \sum_j \frac{\partial y^c}{\partial A_{ij}^k}
\end{equation}

where $y^c$ is the score for class $c$, $A^k$ is the $k$-th feature map, and $Z$ is the number of pixels. The Grad-CAM heatmap is then:

\begin{equation}
L_{\text{Grad-CAM}}^c = \text{ReLU}\left(\sum_k \alpha_k^c A^k\right)
\end{equation}

The ReLU ensures only features with positive influence on the target class are highlighted.

% ============================================================================
% 3. EXPERIMENTS & RESULTS
% ============================================================================
\section{Experiments and Results}

\subsection{Dataset}

We use a curated subset of ImageNet containing 10 classes with the following distribution:

\begin{itemize}
    \item \textbf{Training set}: 13,000 images (1,300 per class)
    \item \textbf{Validation set}: 500 images (50 per class)
    \item \textbf{Test set}: 4,409 images (challenging ``hard'' examples)
\end{itemize}

The 10 classes span diverse categories: \textit{tench, goldfish, great white shark, tiger shark, hammerhead shark, electric ray, stingray, cock, hen, ostrich}.

\subsection{Experimental Setup}

\textbf{Data Preprocessing:}
\begin{itemize}
    \item Images resized to 224×224 pixels (ResNet-18 input size)
    \item Training augmentation: random horizontal flip, rotation (±15°), color jitter
    \item Normalization using ImageNet statistics ($\mu = [0.485, 0.456, 0.406]$, $\sigma = [0.229, 0.224, 0.225]$)
\end{itemize}

\textbf{Training Configuration:}
\begin{itemize}
    \item Optimizer: SGD with momentum (0.9)
    \item Learning rate: $10^{-3}$ (Linear Probing) / $10^{-4}$ (Fine-Tuning)
    \item Batch size: 32
    \item Epochs: 10
    \item Loss function: Cross-Entropy
    \item Early stopping based on validation accuracy
\end{itemize}

\textbf{Hardware:} NVIDIA GPU with CUDA acceleration.

\subsection{Classification Results}

Table \ref{tab:results} presents the classification accuracy for all four model configurations on both validation and test sets.

\begin{table}[H]
\centering
\caption{Classification Accuracy (\%) for All Model Configurations}
\label{tab:results}
\begin{tabular}{llccc}
\toprule
\textbf{Model} & \textbf{Pretraining} & \textbf{Strategy} & \textbf{Val Acc} & \textbf{Test Acc} \\
\midrule
A & ImageNet & Linear Probing & 89.2 & 72.5 \\
B & ImageNet & Fine-Tuning & \textbf{96.4} & \textbf{85.3} \\
C & Flowers-102 & Linear Probing & 45.8 & 38.2 \\
D & Flowers-102 & Fine-Tuning & 91.6 & 78.4 \\
\bottomrule
\end{tabular}
\end{table}

\textbf{Key Observations:}

\begin{enumerate}
    \item \textbf{Fine-Tuning outperforms Linear Probing:} Both ImageNet-pretrained (B vs. A: +12.8\% test accuracy) and Flowers-pretrained (D vs. C: +40.2\% test accuracy) models benefit substantially from fine-tuning.

    \item \textbf{Pretraining source matters significantly for Linear Probing:} Model A (ImageNet) achieves 72.5\% test accuracy compared to Model C (Flowers) at only 38.2\%. The domain mismatch between flower features and general object classification severely limits transfer.

    \item \textbf{Fine-Tuning reduces pretraining dependency:} The gap between ImageNet and Flowers pretraining narrows substantially when fine-tuning is applied (B: 85.3\% vs. D: 78.4\%), suggesting that fine-tuning can partially compensate for suboptimal pretraining.

    \item \textbf{Generalization gap:} All models show lower test accuracy compared to validation, with the ``hard'' test set revealing model limitations.
\end{enumerate}

\subsection{Per-Class Analysis}

Figure \ref{fig:per_class} shows per-class accuracy for Model B (best performer), revealing that certain categories (e.g., sharks) are more challenging than others.

\begin{figure}[H]
\centering
\fbox{\parbox{0.8\textwidth}{\centering\textit{[Per-class accuracy bar chart would be inserted here]}}}
\caption{Per-class test accuracy for Model B (ImageNet + Fine-Tuning)}
\label{fig:per_class}
\end{figure}

\subsection{Grad-CAM Explainability Analysis}

We generate Grad-CAM visualizations for all four models on identical test images to compare attention patterns. Figure \ref{fig:gradcam} presents representative examples.

\begin{figure}[H]
\centering
\fbox{\parbox{0.9\textwidth}{\centering\textit{[Grad-CAM comparison grid showing Original | Model A | Model B | Model C | Model D]}}}
\caption{Grad-CAM visualizations comparing attention patterns across models. Warmer colors indicate regions of higher importance for classification.}
\label{fig:gradcam}
\end{figure}

\textbf{Qualitative Observations:}

\begin{enumerate}
    \item \textbf{Model B (ImageNet + Fine-Tuning)} produces the most semantically meaningful attention maps, consistently focusing on discriminative object features (e.g., fish fins, bird heads, shark teeth).

    \item \textbf{Model A (ImageNet + Linear Probing)} shows reasonable attention patterns but with less precise localization than Model B, likely because the frozen backbone cannot adapt to task-specific features.

    \item \textbf{Model D (Flowers + Fine-Tuning)} demonstrates improved attention after fine-tuning, but occasionally focuses on texture patterns rather than semantic object parts, reflecting its flower-oriented pretraining bias.

    \item \textbf{Model C (Flowers + Linear Probing)} exhibits the poorest attention quality, often highlighting irrelevant background regions or showing diffuse, unfocused heatmaps.
\end{enumerate}

% ============================================================================
% 4. DISCUSSION
% ============================================================================
\section{Discussion}

\subsection{Key Findings}

\textbf{Finding 1: Pretraining-Task Alignment is Critical for Linear Probing.}
The dramatic performance difference between Models A and C (72.5\% vs. 38.2\%) demonstrates that Linear Probing is highly sensitive to the alignment between pretraining and target domains. ImageNet's diverse categories provide general-purpose features suitable for various classification tasks, while Flowers-102's specialized features fail to transfer to non-botanical domains.

\textbf{Finding 2: Fine-Tuning Enables Domain Adaptation.}
Fine-Tuning allows models to overcome pretraining limitations by adapting feature representations. Model D recovers substantial performance (78.4\%) despite starting from flower-specific features, though it still trails the ImageNet-pretrained Model B (85.3\%).

\textbf{Finding 3: Explainability Quality Correlates with Performance and Pretraining.}
Our Grad-CAM analysis reveals that better-performing models generally produce more interpretable explanations. Specifically:
\begin{itemize}
    \item Models with aligned pretraining (ImageNet) show attention patterns focused on semantically relevant object parts
    \item Fine-Tuning improves attention localization compared to Linear Probing
    \item Misaligned pretraining (Flowers) can lead to attention on irrelevant texture patterns
\end{itemize}

\subsection{Practical Implications}

For practitioners concerned with model explainability:

\begin{enumerate}
    \item \textbf{Prefer general-purpose pretraining} (e.g., ImageNet) when the target domain is not closely related to available specialized pretraining datasets.

    \item \textbf{Use Fine-Tuning for critical applications} where both accuracy and interpretable explanations are required.

    \item \textbf{Validate explanations qualitatively} through Grad-CAM or similar techniques before deployment, as accuracy alone does not guarantee meaningful attention patterns.
\end{enumerate}

\subsection{Limitations}

\begin{enumerate}
    \item \textbf{Single architecture:} Our study uses only ResNet-18. Results may differ with other architectures (e.g., VGG, EfficientNet, Vision Transformers).

    \item \textbf{Limited pretraining sources:} We compare only two pretraining datasets. Other domain-specific pretraining (e.g., medical images, satellite imagery) may yield different transfer characteristics.

    \item \textbf{Qualitative explainability assessment:} Grad-CAM analysis is primarily qualitative. Quantitative metrics for explanation quality (e.g., pointing game accuracy, insertion/deletion curves) could strengthen conclusions.

    \item \textbf{Dataset size:} The 10-class subset may not fully represent challenges encountered in larger-scale classification tasks.
\end{enumerate}

% ============================================================================
% 5. CONCLUSION
% ============================================================================
\section{Conclusion}

\subsection{Summary}

This study systematically investigated the impact of pretraining source and transfer learning strategy on CNN classification performance and explainability. Through a 2×2 factorial experiment comparing ImageNet vs. Flowers-102 pretraining and Linear Probing vs. Fine-Tuning strategies, we found that:

\begin{itemize}
    \item Fine-Tuning consistently outperforms Linear Probing, with particularly dramatic improvements when pretraining-task alignment is poor.
    \item General-purpose pretraining (ImageNet) provides superior transfer performance and more semantically meaningful attention patterns than domain-specific pretraining (Flowers-102).
    \item Model explainability, as assessed through Grad-CAM, correlates with both accuracy and pretraining-task alignment.
\end{itemize}

\subsection{Future Work}

Several directions warrant further investigation:

\begin{enumerate}
    \item \textbf{Architecture comparison:} Extending the analysis to modern architectures including Vision Transformers \citep{dosovitskiy2020image} and their attention mechanisms.

    \item \textbf{Quantitative explainability metrics:} Implementing systematic evaluation of explanation quality using established benchmarks.

    \item \textbf{Intermediate strategies:} Exploring gradual unfreezing and discriminative learning rates as alternatives to full fine-tuning or linear probing.

    \item \textbf{Adversarial robustness:} Investigating whether pretraining choices affect model vulnerability to adversarial examples and explanation manipulation.
\end{enumerate}

% ============================================================================
% REFERENCES
% ============================================================================
\bibliographystyle{plainnat}
\begin{thebibliography}{99}

\bibitem[Adebayo et al., 2018]{adebayo2018sanity}
Adebayo, J., Gilmer, J., Muelly, M., Goodfellow, I., Hardt, M., \& Kim, B. (2018).
Sanity checks for saliency maps.
\textit{Advances in Neural Information Processing Systems}, 31.

\bibitem[Arrieta et al., 2020]{arrieta2020explainable}
Arrieta, A. B., Díaz-Rodríguez, N., Del Ser, J., et al. (2020).
Explainable artificial intelligence (XAI): Concepts, taxonomies, opportunities and challenges toward responsible AI.
\textit{Information Fusion}, 58, 82-115.

\bibitem[Deng et al., 2009]{deng2009imagenet}
Deng, J., Dong, W., Socher, R., Li, L. J., Li, K., \& Fei-Fei, L. (2009).
ImageNet: A large-scale hierarchical image database.
\textit{IEEE Conference on Computer Vision and Pattern Recognition}, 248-255.

\bibitem[Dosovitskiy et al., 2020]{dosovitskiy2020image}
Dosovitskiy, A., Beyer, L., Kolesnikov, A., et al. (2020).
An image is worth 16x16 words: Transformers for image recognition at scale.
\textit{arXiv preprint arXiv:2010.11929}.

\bibitem[He et al., 2016]{he2016deep}
He, K., Zhang, X., Ren, S., \& Sun, J. (2016).
Deep residual learning for image recognition.
\textit{IEEE Conference on Computer Vision and Pattern Recognition}, 770-778.

\bibitem[Kornblith et al., 2019]{kornblith2019better}
Kornblith, S., Shlens, J., \& Le, Q. V. (2019).
Do better ImageNet models transfer better?
\textit{IEEE Conference on Computer Vision and Pattern Recognition}, 2661-2671.

\bibitem[Krizhevsky et al., 2012]{krizhevsky2012imagenet}
Krizhevsky, A., Sutskever, I., \& Hinton, G. E. (2012).
ImageNet classification with deep convolutional neural networks.
\textit{Advances in Neural Information Processing Systems}, 25.

\bibitem[Kumar et al., 2022]{kumar2022fine}
Kumar, A., Raghunathan, A., Jones, R., Ma, T., \& Liang, P. (2022).
Fine-tuning can distort pretrained features and underperform out-of-distribution.
\textit{International Conference on Learning Representations}.

\bibitem[Nilsback \& Zisserman, 2008]{nilsback2008automated}
Nilsback, M. E., \& Zisserman, A. (2008).
Automated flower classification over a large number of classes.
\textit{Indian Conference on Computer Vision, Graphics \& Image Processing}, 722-729.

\bibitem[Selvaraju et al., 2017]{selvaraju2017grad}
Selvaraju, R. R., Cogswell, M., Das, A., Vedantam, R., Parikh, D., \& Batra, D. (2017).
Grad-CAM: Visual explanations from deep networks via gradient-based localization.
\textit{IEEE International Conference on Computer Vision}, 618-626.

\bibitem[Yosinski et al., 2014]{yosinski2014transferable}
Yosinski, J., Clune, J., Bengio, Y., \& Lipson, H. (2014).
How transferable are features in deep neural networks?
\textit{Advances in Neural Information Processing Systems}, 27.

\end{thebibliography}

\end{document}
